A \emph{transformation} of degree $n\in \N$ is a function from $\{1, \ldots,
n\}$ to itself. A \emph{transformation semigroup} is a semigroup consisting
of transformations of equal degree and with the operation of composition of
functions. For the sake of brevity we will denote $\{1,\ldots,n\}$ by $\nset$.
We define $(a \to b)$ to be the transformation defined by
\begin{equation*} 
v (a \to b) = 
 \begin{cases}
 b & \text{if } v = a\\
 v & \text{otherwise}
 \end{cases} 
\end{equation*}
where $a, b \in \nset$ and $a \ne b$. A \emph{digraph} is an ordered pair
$(V, A)$, where $V$ is a set whose elements are referred to as
\emph{vertices}, and $A\subseteq (V\times V) \setminus \{(v,v) : v \in V\}$
is a set of ordered pairs called
\emph{arcs}. We identify a transformation $(a \to b)$ with an arc $(a, b)$
in a digraph and we refer to $(a \to b)$ as an arc. If $D$ is a digraph,
we denote by $\langle D \rangle$ the semigroup generated by
the arcs of $D$, and we refer to such a semigroup as \emph{arc-generated}.

If $A$ is empty, then $D$ is an empty digraph, and $\genset{D}$ is the empty
semigroup. Because this is a degenerate case, and in order to simplify the
statement of our results, we shall always assume that $D$ is not empty (and
hence, that $\genset{D}$ is not the empty semigroup). 